% Gerado por regressao/06_tabelas_latex.R (projeto Modelagem). Não editar à mão.
\begin{landscape}
\begin{table}[htbp]
\caption{Robustez R3 — variável dependente em log}\label{tab:robustez_R3}
\centering
\footnotesize
\setlength{\tabcolsep}{3pt}
\begin{tabular}{@{}>{\raggedright\arraybackslash}p{\dimexpr\linewidth-594pt\relax}>{\centering\arraybackslash}p{60pt}>{\centering\arraybackslash}p{60pt}>{\centering\arraybackslash}p{60pt}>{\centering\arraybackslash}p{60pt}>{\centering\arraybackslash}p{60pt}>{\centering\arraybackslash}p{60pt}>{\centering\arraybackslash}p{60pt}>{\centering\arraybackslash}p{60pt}>{\centering\arraybackslash}p{60pt}@{}}
\hline
\textbf{Variável} & \textbf{\hspace{0pt}Despesa total} & \textbf{\hspace{0pt}Saúde e Saneamento} & \textbf{\hspace{0pt}Educação e Cultura} & \textbf{\hspace{0pt}Habitação e Urbanismo} & \textbf{\hspace{0pt}Assistência Social} & \textbf{\hspace{0pt}Transporte} & \textbf{\hspace{0pt}Administração} & \textbf{\hspace{0pt}Agricultura} & \textbf{\hspace{0pt}Comunicações} \\ \hline
\multicolumn{10}{@{}l}{\textit{Modelo A (ciclo eleitoral)}} \\
Ano eleitoral & 0,054\textsuperscript{***} & 0,043\textsuperscript{***} & 0,038\textsuperscript{**} & 0,241\textsuperscript{***} & 0,071\textsuperscript{***} & 0,199\textsuperscript{***} & 0,010 & 0,061\textsuperscript{***} & 0,071\textsuperscript{***} \\
 & (0,011) & (0,004) & (0,016) & (0,020) & (0,012) & (0,025) & (0,009) & (0,017) & (0,015) \\
Ano pré-eleitoral & 0,023 & $-$0,007 & 0,048\textsuperscript{**} & 0,055\textsuperscript{**} & 0,037\textsuperscript{*} & 0,023 & 0,009 & 0,066\textsuperscript{***} & 0,079\textsuperscript{***} \\
 & (0,015) & (0,007) & (0,023) & (0,024) & (0,020) & (0,028) & (0,009) & (0,024) & (0,020) \\
\hline
Observações & 70.669 & 70.526 & 70.550 & 69.574 & 70.484 & 54.421 & 70.553 & 64.463 & 12.358 \\
Municípios & 5.568 & 5.568 & 5.568 & 5.564 & 5.568 & 5.205 & 5.568 & 5.435 & 1.894 \\
R\textsuperscript{2} & 0,840 & 0,665 & 0,656 & 0,242 & 0,368 & 0,099 & 0,299 & 0,102 & 0,020 \\[6pt]
\multicolumn{10}{@{}l}{\textit{Modelo B (coalizões)}} \\
Prefeito na coalizão & 0,003\textsuperscript{***} & 0,001 & 0,003\textsuperscript{*} & 0,016\textsuperscript{**} & 0,000 & 0,011 & 0,006\textsuperscript{**} & 0,012 & 0,006 \\
do governador & (0,001) & (0,002) & (0,002) & (0,007) & (0,003) & (0,012) & (0,003) & (0,009) & (0,029) \\
Prefeito na coalizão & 0,002 & 0,003 & 0,000 & $-$0,012 & $-$0,004 & 0,054\textsuperscript{***} & 0,007\textsuperscript{*} & 0,012 & $-$0,013 \\
do presidente & (0,001) & (0,003) & (0,002) & (0,009) & (0,004) & (0,017) & (0,004) & (0,012) & (0,042) \\
\hline
Observações & 70.669 & 70.526 & 70.550 & 69.574 & 70.484 & 54.421 & 70.553 & 64.463 & 12.358 \\
Municípios & 5.568 & 5.568 & 5.568 & 5.564 & 5.568 & 5.205 & 5.568 & 5.435 & 1.894 \\
R\textsuperscript{2} & 0,344 & 0,140 & 0,153 & 0,029 & 0,061 & 0,009 & 0,072 & 0,014 & 0,009 \\
\hline
\end{tabular}
\fonte{Elaboração própria.}
\nota{Modelo A: efeito fixo de município, erros-padrão de Driscoll-Kraay. Modelo B: efeitos fixos de município e de ano, erros-padrão agrupados por município. Erros-padrão entre parênteses. \textsuperscript{***}p$<$0,01; \textsuperscript{**}p$<$0,05; \textsuperscript{*}p$<$0,1.}
\end{table}
\end{landscape}
